% Autogenerated translation of PROJECT_GUIDE_zh.md by Texpad
% To stop this file being overwritten during the typeset process, please move or remove this header

\documentclass[12pt]{book}
\usepackage{graphicx}
\usepackage{fontspec}
\usepackage[utf8]{inputenc}
\usepackage[a4paper,left=.5in,right=.5in,top=.3in,bottom=0.3in]{geometry}
\setlength\parindent{0pt}
\setlength{\parskip}{\baselineskip}
\setmainfont{Helvetica Neue}
\usepackage{hyperref}
\pagestyle{plain}
\begin{document}

\chapter*{UnifoLM-WLA 源码导读（导师笔记）}

\begin{quote}
面向"从零理解"的读者。本文按 \textbf{数据层 → 模型层 → 训练层} 的顺序拆解仓库里的每一个文件，
给出端到端的数据形状追踪、关键设计取舍的原因，以及我在通读中发现的问题与死代码清单。
所有引用格式为 \texttt{文件:行号}。
\end{quote}

\hrule
\section*{目录}

\begin{itemize}
\item \href{#0-一句话定位}{0. 一句话定位}
\item \href{#1-全局地图}{1. 全局地图}
\item \href{#2-数据层}{2. 数据层}
\item \href{#3-模型层}{3. 模型层}
\item \href{#4-训练层}{4. 训练层}
\item \href{#5-端到端-trace手算维度}{5. 端到端 Trace（手算维度）}
\item \href{#6-关键设计决策为什么这么做}{6. 关键设计决策：为什么这么做}
\item \href{#7-坑与风险清单}{7. 坑与风险清单}
\item \href{#8-死代码--未接线清单}{8. 死代码 / 未接线清单}
\item \href{#9-上手路线图}{9. 上手路线图}
\item \href{#10-自测题}{10. 自测题}
\end{itemize}

\hrule
\section*{0. 一句话定位}

UnifoLM-WLA 是宇树（Unitree）的人形机器人 VLA（Vision-Language-Action）基础模型训练代码库：
\textbf{用 Qwen3-VL 当"眼睛+语言理解"，用 MMDiT 双流联合注意力的 flow-matching 头当"运动皮层"，预测未来 1 秒的连续动作块。}

\begin{itemize}
\item 论文/模型卡：6B 参数，\textasciitilde{}2500 小时真机数据，64 个任务，支持二指夹爪与五指灵巧手（README.md:15）
\item 本仓库\textbf{只开源了动作专家（action expert）的训练代码}，不是完整推理部署栈
\item 底座借用 \href{https://github.com/starVLA/starVLA}{starVLA} 的分层思路与 \href{https://github.com/QwenLM/Qwen-Image}{Qwen-Image} 的 MMDiT 实现（README.md:54-57）
\end{itemize}

\textbf{核心范式}：VLM 是\textbf{冻结/微调的编码器}（不是自回归生成器），它最后一层的 hidden states 作为 DiT 的
cross-attention 条件（key/value）；DiT 以纯噪声为起点，用流匹配（flow matching）在 4 步内积分出动作轨迹。

\hrule
\section*{1. 全局地图}

\texttt{
                       ┌──────────────────────────────────────────┐
   LeRobot 数据集 ───▶ │ 数据层 dataloader/multi\_source\_dataset/   │
   (parquet + mp4)     │  统一 54 维动作 / 60 维状态 / 掩码 / 图像  │
                       └───────────────────┬──────────────────────┘
                                           │ collate → dict batch
                                           ▼
   ┌───────────────────────────────────────────────────────────────────┐
   │ 模型层 model/                                                      │
   │  ┌─────────────────────┐        ┌──────────────────────────────┐  │
   │  │ QWen3.py            │        │ MMDiT\_ActionHeader.py        │  │
   │  │ Qwen3-VL 编码器     │        │  ActionEncoder + DiT + Decoder│  │
   │  │ (+ robot\_state 投影)│        │  flow matching loss          │  │
   │  └─────────┬───────────┘        └───────────▲──────────────────┘  │
   │            │ last\_hidden\_states              │ cross-attn K/V     │
   │            └─────────────────────────────────┘                    │
   │  组装者: framework/VLM4A/QwenMMDiT.py                              │
   └───────────────────────────────┬───────────────────────────────────┘
                                   ▼
   ┌───────────────────────────────────────────────────────────────────┐
   │ 训练层 training/                                                   │
   │  train\_unifolm\_wla.py: VLATrainer 主循环                           │
   │  trainer\_utils/: Accelerator+DeepSpeed、config 访问追踪、日志       │
   └───────────────────────────────────────────────────────────────────┘
}

\subsection*{文件清单与职责（36 个 .py，共 \textasciitilde{}7700 行）}

$\vert$ 层 $\vert$ 文件 $\vert$ 行数 $\vert$ 职责 $\vert$
$\vert$---$\vert$---$\vert$---$\vert$---$\vert$
$\vert$ \textbf{入口} $\vert$ \texttt{training/train\_unifolm\_wla.py} $\vert$ 508 $\vert$ \texttt{main()} + \texttt{VLATrainer} 训练循环 $\vert$
$\vert$ \textbf{模型组装} $\vert$ \texttt{model/framework/base\_framework.py} $\vert$ 269 $\vert$ \texttt{build\_framework()} 注册表工厂 + \texttt{baseframework} 基类 $\vert$
$\vert$ $\vert$ \texttt{model/framework/VLM4A/QwenMMDiT.py} $\vert$ 386 $\vert$ \textbf{唯一注册的框架}：拼装 VLM + 动作头，定义 forward/predict $\vert$
$\vert$ $\vert$ \texttt{model/framework/share\_tools.py} $\vert$ 559 $\vert$ 配置合并、checkpoint 读取、config 版本兼容 $\vert$
$\vert$ $\vert$ \texttt{model/tools.py} $\vert$ 428 $\vert$ \texttt{Registry}、模块冻结工具、\texttt{CrossAttention}（死代码） $\vert$
$\vert$ \textbf{VLM} $\vert$ \texttt{model/modules/vlm/QWen3.py} $\vert$ 369 $\vert$ Qwen3-VL 封装 + \textbf{提示词模板} + robot\_state 注入 hook $\vert$
$\vert$ $\vert$ \texttt{model/modules/vlm/\_\_init\_\_.py} $\vert$ 10 $\vert$ 按模型名分派（只支持 Qwen3-VL） $\vert$
$\vert$ $\vert$ \texttt{model/modules/projector/robot\_state.py} $\vert$ 239 $\vert$ 状态→VLM embedding 的 MLP 投影器 $\vert$
$\vert$ \textbf{动作头} $\vert$ \texttt{model/modules/action\_model/MMDiT\_ActionHeader.py} $\vert$ 367 $\vert$ flow matching 的编码/解码/损失/采样 $\vert$
$\vert$ $\vert$ \texttt{model/modules/action\_model/DiT\_modules/mmdit.py} $\vert$ 593 $\vert$ DiT 主干 + 双流联合注意力 + RoPE $\vert$
$\vert$ $\vert$ \texttt{.../flow\_matching\_head/action\_encoder.py} $\vert$ 84 $\vert$ 正弦时间编码 + swish $\vert$
$\vert$ $\vert$ \texttt{.../flow\_matching\_head/cross\_attention\_dit.py} $\vert$ 362 $\vert$ \textbf{完全未被引用的旧实现} $\vert$
$\vert$ \textbf{数据} $\vert$ \texttt{dataloader/\_\_init\_\_.py} $\vert$ 50 $\vert$ \texttt{build\_dataloader()} 单一入口 + 落盘统计量 $\vert$
$\vert$ $\vert$ \texttt{.../single\_source\_dataset.py} $\vert$ 750 $\vert$ \textbf{数据层核心}：单数据源 → 统一表示 $\vert$
$\vert$ $\vert$ \texttt{.../dataloader.py} $\vert$ 235 $\vert$ \texttt{collate\_fn} + DataLoader 构造 $\vert$
$\vert$ $\vert$ \texttt{.../action\_mapping.py} $\vert$ 253 $\vert$ 统一动作/状态槽位定义与填充 $\vert$
$\vert$ $\vert$ \texttt{.../se3\_utils.py} $\vert$ 180 $\vert$ 位姿格式互转 + SE(3) 相对变换 $\vert$
$\vert$ $\vert$ \texttt{.../stats\_utils.py} $\vert$ 102 $\vert$ 统计量加载 + 三种归一化器 $\vert$
$\vert$ $\vert$ \texttt{.../lerobot\_wrapper.py} $\vert$ 235 $\vert$ LeRobotDataset 选择性加载 + Arrow 缓存 $\vert$
$\vert$ $\vert$ \texttt{.../config.py} $\vert$ 226 $\vert$ 数据 YAML → dataclass $\vert$
$\vert$ $\vert$ \texttt{.../distributed\_utils.py} $\vert$ 100 $\vert$ 多卡缓存构建串行化 $\vert$
$\vert$ $\vert$ \texttt{.../fps\_utils.py} $\vert$ 83 $\vert$ \textbf{完全未被引用的帧率重采样} $\vert$
$\vert$ \textbf{训练工具} $\vert$ \texttt{training/trainer\_utils/trainer\_tools.py} $\vert$ 558 $\vert$ \texttt{TrainerUtils}（冻结/加载/打印）+ 一堆死代码 $\vert$
$\vert$ $\vert$ \texttt{training/trainer\_utils/config\_tracker.py} $\vert$ 548 $\vert$ \texttt{AccessTrackedConfig} + \textbf{monkey patch OmegaConf} $\vert$
$\vert$ $\vert$ \texttt{training/trainer\_utils/overwatch.py} $\vert$ 168 $\vert$ 日志 + 关键的 DataLoader worker 防 CUDA 检测 $\vert$

\hrule
\section*{2. 数据层}

\subsection*{2.1 数据来源：LeRobot 格式}

\texttt{
$<$data\_base$>$/$<$data\_path$>$/
├── meta/info.json          \# features 定义（含图像分辨率）
├── data/chunk-000/*.parquet
└── videos/...
}

训练数据的 YAML 在 \texttt{unifolm\_wla/dataloader/multi\_source\_dataset/configs/unitree.yaml}。
\texttt{multi\_task: true} 时，\texttt{\_load\_datasets()} (\texttt{single\_source\_dataset.py:139}) 用
\texttt{discover\_task\_dirs()} 扫描 \texttt{data\_path} 下所有含 \texttt{meta/info.json} 的子目录，每个任务一个 \texttt{LeRobotDataset}，
再用 \texttt{\_route\_index()} (\texttt{single\_source\_dataset.py:436}) 把全局 index 线性映射到 (子数据集, 局部 index)。

\texttt{dataloader.py:171-175} 再把所有 \texttt{enabled: true} 的数据源用 \texttt{ConcatDataset} 拼接。
\textbf{注意}：采样比例 = 各源数据量占比，配置里的 \texttt{weight} 字段当前完全不生效（\texttt{dataloader.py:141-142} 有注释说明）。

\subsection*{2.2 未来动作块怎么来的}

\texttt{\_build\_delta\_timestamps()} (\texttt{single\_source\_dataset.py:121})：

\texttt{python
future\_ts = [i / fps for i in range(0, fps)]   \# fps=30 → 30 个时间戳，覆盖 1 秒
dt[action\_key] = future\_ts                      \# 只对 action key 生效
}

LeRobot 的 \texttt{delta\_timestamps} 机制会把该特征展开成 \texttt{(T, D)}，即"当前帧 + 未来 29 帧"。
状态 key 不加 delta → 仍是当前帧的单个向量。

\begin{quote}
⚠️ \texttt{chunk\_size}（配置里 30）\textbf{从未被用来切片}，实际长度恒等于 \texttt{source\_fps}。
\end{quote}

\subsection*{2.3 统一动作空间：54 维（\texttt{action\_mapping.py:8-43}）}

$\vert$ 槽位 $\vert$ 维度 $\vert$ 含义 $\vert$ 处理方式 $\vert$
$\vert$---$\vert$---$\vert$---$\vert$---$\vert$
$\vert$ \texttt{[0:6]} $\vert$ 6 $\vert$ 左臂 EE（xyz + rotvec） $\vert$ \textbf{相对当前状态的 SE(3) 变换} $\vert$
$\vert$ \texttt{[6:7]} $\vert$ 1 $\vert$ 左夹爪 $\vert$ 绝对，归一化 $\vert$
$\vert$ \texttt{[7:13]} $\vert$ 6 $\vert$ 左手灵巧指 6-DoF $\vert$ 透传 $\vert$
$\vert$ \texttt{[13:19]} $\vert$ 6 $\vert$ 右臂 EE $\vert$ 相对 SE(3) $\vert$
$\vert$ \texttt{[19:20]} $\vert$ 1 $\vert$ 右夹爪 $\vert$ 绝对 $\vert$
$\vert$ \texttt{[20:26]} $\vert$ 6 $\vert$ 右手灵巧指 6-DoF $\vert$ 透传 $\vert$
$\vert$ \texttt{[26:29]} $\vert$ 3 $\vert$ waist 关节 $\vert$ 绝对 $\vert$
$\vert$ \texttt{[29:32]} $\vert$ 3 $\vert$ torso 关节 $\vert$ 绝对 $\vert$
$\vert$ \texttt{[32:34]} $\vert$ 2 $\vert$ 底盘 vx, vy $\vert$ 绝对 $\vert$
$\vert$ \texttt{[34:35]} $\vert$ 1 $\vert$ 底盘 vw（角速度） $\vert$ 绝对 $\vert$
$\vert$ \texttt{[35:41]} $\vert$ 6 $\vert$ 底盘位姿相对变换 $\vert$ \textbf{相对 SE(3)} $\vert$
$\vert$ \texttt{[41:42]} $\vert$ 1 $\vert$ 身高指令 $\vert$ 绝对 $\vert$
$\vert$ \texttt{[42:48]} $\vert$ 6 $\vert$ 左腿关节 $\vert$ 绝对 $\vert$
$\vert$ \texttt{[48:54]} $\vert$ 6 $\vert$ 右腿关节 $\vert$ 绝对 $\vert$

外加一个 \textbf{54 维 bool 掩码}：只标记该数据源真实存在的模块；不存在的槽位恒为 0 且掩码为 False。
这个掩码一路传到 loss 里当权重用（\texttt{MMDiT\_ActionHeader.py:285-290}），使不同本体的数据能安全混合训练。

\textbf{单臂数据复用左臂槽位}（\texttt{action\_mapping.py:129-138}）——这带来一个隐患，见 §6.3。

\subsection*{2.4 统一状态空间：60 维（\texttt{action\_mapping.py:47-62}）}

与动作槽位语义一一对应，但有两处关键差异：

\begin{enumerate}
\item \textbf{EE 姿态用 rotation-6D（9 维 = xyz + 6）而非 rotvec（6 维）}
（\texttt{STATE\_SLICES["left\_xyz\_rot6d"] = slice(0, 9)}）
原因：rot6d 无奇异点、连续，更适合当网络输入；而动作输出用 rotvec 更紧凑、更接近增量语义。
\texttt{matrix\_to\_rot6d()} (\texttt{se3\_utils.py:65}) 取旋转矩阵前两列并转置成 \texttt{[R00,R10,R20,R01,R11,R21]}。
\item \texttt{[41:47]} 的 \texttt{base\_rotvec} 在状态里实际存的是
\textbf{重力方向 xyz(3) + 机体角速度 xyz(3)}（\texttt{action\_mapping.py:239-243}），不是位姿。
\end{enumerate}

\begin{quote}
📌 同一个 \texttt{[48:54]} 槽位，在\textbf{动作向量里是右腿}，在\textbf{状态向量里是左腿}（状态前面每个模块多占 3 维）。
读代码时务必先确认是哪个向量。文档 \texttt{docs/robot\_action\_state\_processing.md:198-202} 专门强调了这点。
\end{quote}

\subsection*{2.5 相对动作 = SE(3) 变换（\texttt{se3\_utils.py:162-180}）}

这是整个数据层最核心的数学：

\texttt{python
T\_curr  = pose\_to\_se3\_from\_format(current\_state\_pose, ee\_format)   \# (4,4)
T\_future= pose\_to\_se3\_from\_format(future\_action\_poses, ee\_format)  \# (N,4,4)
T\_rel   = se3\_inverse(T\_curr) @ T\_future                           \# (N,4,4)
return se3\_to\_xyz\_rotvec(T\_rel)                                    \# (N,6)
}

即 \texttt{T\_rel = T\_curr⁻¹ · T\_future}。\textbf{为什么用相对？}
绝对位姿强依赖机器人基座标系，跨本体/跨场地不可迁移；相对位姿是"从当前位置往哪走"，天然与坐标系解耦，
是跨本体混合训练的前提。\texttt{se3\_inverse} 用解析式 \texttt{[[Rᵀ, -Rᵀt],[0,1]]} 而不是求逆矩阵，更快更稳。

源数据的 \texttt{ee\_format} 由配置声明（\texttt{xyz\_rpy} / \texttt{xyz\_quat} / \texttt{xyz\_rvec}），Unitree 是 \texttt{xyz\_rpy}。

\subsection*{2.6 归一化（\texttt{stats\_utils.py:38-97}）}

离线预先算好两个 json，放在 \texttt{dataloader/multi\_source\_dataset/stats/}：

\begin{itemize}
\item \textbf{\texttt{stats.json}}：普通特征（关节、夹爪、body command）的 min/max/mean/std/q01/q99
\item **\texttt{relative\_stats.json}**：相对位姿动作的统计量，键为 \texttt{left\_ee\_pose\_gripper\_base} 等
\end{itemize}

支持三种归一化器，输出统一为 \texttt{normalized = (value - offset) / scale}：

$\vert$ 类型 $\vert$ offset $\vert$ scale $\vert$ 用于 $\vert$
$\vert$---$\vert$---$\vert$---$\vert$---$\vert$
$\vert$ \texttt{minmax\_q} $\vert$ \texttt{(q01+q99)/2} $\vert$ \texttt{(q99-q01)/2} $\vert$ \texttt{norm\_type} / \texttt{state\_norm\_type} / \texttt{gripper\_norm\_type} $\vert$
$\vert$ \texttt{zscore} $\vert$ \texttt{mean} $\vert$ \texttt{std} $\vert$ \texttt{rel\_norm\_type}（相对动作默认） $\vert$
$\vert$ \texttt{minmax} $\vert$ \texttt{(min+max)/2} $\vert$ \texttt{(max-min)/2} $\vert$ 备选 $\vert$

\texttt{get\_normalizer} 优先取 \texttt{global\_*}（跨任务合并后的全局量，shape=(6,)），回退到 per-step 的
\texttt{q01/q99}（shape=(30,6)）。\texttt{relative\_stats.json} 两者都有，所以走 \texttt{global\_*}；
这正好对应文档 §7「同一本体不同任务统计量按等任务权重合并」。

\textbf{两处特判值得记}：

\begin{itemize}
\item \texttt{base\_pose} 的状态归一化被强制改成 \texttt{offset=0, scale=1}（\texttt{single\_source\_dataset.py:340-342}，带 \texttt{(TODO) HACK} 标注）
\item \texttt{base\_rot} 只归一化后 3 维角速度，重力方向保持原样（\texttt{single\_source\_dataset.py:344-352}）
\end{itemize}

\textbf{关于权重的三个"必须"}：
\texttt{\_action\_norm\_offset/\_scale} 是 54 维向量，按槽位把各模块的 offset/scale 拼起来
（\texttt{\_build\_action\_norm\_vectors()}），本意是让每个样本自带反归一化参数。但见 §8：它被算出来、被 collate，
\textbf{却没有消费者}。

\subsection*{2.7 collate：从样本列表到训练张量（\texttt{dataloader.py:45-124}）}

\texttt{collate\_fn} 是理解 batch 结构的关键：

\texttt{python
\{
  "action":      (B, T, 54)      \# T = source\_fps = 30
  "action\_mask": (B, 54)
  "state":       (B, 60)
  "state\_mask":  (B, 60)
  "action\_norm\_offset": (B, 54)  \# ← 无消费者
  "action\_norm\_scale":  (B, 54)  \# ← 无消费者
  "images":    \{role: (B, C, H, W)\}   \# 缺失相机用 zeros 填充
  "image\_mask":\{role: (B,)\}           \# bool
  "task":      [str] * B
  "arm\_type":  [str] * B
  "robot\_type":[str] * B
\}
}

两个设计要点：

\begin{enumerate}
\item \textbf{role 顺序必须稳定}：\texttt{all\_roles} 按"首次出现顺序"而非字母序收集（\texttt{dataloader.py:64-75}），
因为 \texttt{\_process\_images()} 是按 \texttt{config.image\_keys} 的顺序建 dict 的，这个顺序必须和 prompt 里
\texttt{role: $<$image$>$} 的顺序严格一致——\textbf{像素顺序错位会让模型看到张冠李戴的画面}。
\item \textbf{同一 role 的 shape 必须一致}，否则直接 \texttt{raise ValueError}（\texttt{dataloader.py:98-102}）。
设计者明确拒绝"静默 padding 混合分辨率"，认为那会浪费算力并给视觉编码器引入黑边伪影。
\end{enumerate}

\subsection*{2.8 两个工程暗坑的处理}

\textbf{(a) Arrow 缓存}（\texttt{lerobot\_wrapper.py:145-198}）
\texttt{LeRobotDatasetWithSelectedKeys} 重写了 \texttt{load\_hf\_dataset()}，把 parquet 读进 pyarrow 后
\texttt{save\_to\_disk} 到 \texttt{cache\_dir/arrow\_cache/$<$md5$>$}，并用 \texttt{.complete} 标记文件做原子提交
（先写 \texttt{.tmp.$<$pid$>$.$<$uuid$>$} 再 \texttt{os.rename}）。为什么要绕这一圈？
注释里写明：\texttt{load\_from\_disk} \textbf{不走 FileLock}，从而躲开 HF datasets 在 CPFS 上的 stale-dentry 竞争。

\textbf{(b) DataLoader worker 里那 520 MiB 显存}（\texttt{dataloader.py:16-42} + \texttt{overwatch.py:146-168}）
这是我在这个仓库里看到的最"血腥"的一段工程注释。问题链条是：

\begin{enumerate}
\item \texttt{multiprocessing\_context="spawn"} 的 worker 会重新执行训练入口脚本
\item 重新 import 时 \texttt{tools.py} 建了模块级 \texttt{overwatch = initialize\_overwatch(\_\_name\_\_)}
\item \texttt{DistributedOverwatch.\_\_init\_\_} 里的 \texttt{accelerate.PartialState()} 触发 \texttt{torch.cuda.set\_device()}
\item \textbf{每个 worker 因此绑定一个 \textasciitilde{}520 MiB 的 CUDA context}，而数据通路是纯 CPU 的
\end{enumerate}

修法很精细：worker 早期\textbf{无法}用 \texttt{sys.argv}（已被父进程覆盖）或
\texttt{type(current\_process())}（\texttt{SpawnProcess} 类型在 re-import 完成后才安装）来判断自己是不是 worker，
唯一可靠信号是 \texttt{current\_process().name == "SpawnProcess-N"}。
\texttt{overwatch.py:165-167} 就用这个判断返回不含 CUDA 的 \texttt{PureOverwatch}。
\texttt{dataloader.py:33-42} 是第二个优化：只在 worker spawn 期间设 \texttt{DS\_ACCELERATOR=cpu}，
跳过 DeepSpeed 的 \texttt{gcc -laio} 兼容性探测（\textbf{不能}用清空 \texttt{CUDA\_VISIBLE\_DEVICES}，会死锁）。

\hrule
\section*{3. 模型层}

\subsection*{3.1 注册表与工厂模式}

\texttt{model/tools.py:117-145} 定义了一个极简 \texttt{Registry}。\texttt{base\_framework.py:26-48} 的
\texttt{\_auto\_import\_framework\_modules()} 会\textbf{动态扫描} \texttt{framework/} 目录并 import 所有模块（跳过 \texttt{\_} 开头的），
从而触发 \texttt{@FRAMEWORK\_REGISTRY.register("QwenMMDiT")} 装饰器（\texttt{QwenMMDiT.py:176}）。

```python

\chapter*{base\_framework.py:51-72}

def build\emph{framework(cfg):        \# 唯一入口
    \_auto}import\emph{framework}modules()
    framework\emph{id = cfg.framework.name          \# ← 必须匹配注册名
    return FRAMEWORK}REGISTRY\href{cfg}{framework\_id}
```

\begin{quote}
⚠️ 当前仓库*\emph{只注册了 \texttt{QwenMMDiT}*}。而 \texttt{config/training/unifolm\_wla\_pretrain\_multisource.yaml:9}
写的是 \texttt{name: QwenGR00T} —— 用这份配置直接跑会在 \texttt{build\_framework} 抛 \texttt{NotImplementedError}。
猜测是开源时从内部代码裁剪、顺手改了类名但漏改 YAML。
\end{quote}

\subsection*{3.2 VLM 编码器：\texttt{QWen3.py}}

\texttt{\_QWen3\_VL\_Interface} 是对 \texttt{Qwen3VLForConditionalGeneration} 的薄封装，做三件事：

\textbf{(1) 加载与对齐}（\texttt{QWen3.py:54-93}）
\texttt{flash\_attention\_2} 若未安装自动回退 \texttt{sdpa}；把 \texttt{config.hidden\_size} 对齐到 \texttt{text\_config.hidden\_size}
（Qwen3-VL 顶层 config 没有 \texttt{hidden\_size} 直通）；\texttt{padding\_side = "right"}。

\textbf{(2) 提示词模板 —— 这是全仓库最该背下来的一段}（\texttt{QWen3.py:216-341}）

单样本 user turn 长这样：

\texttt{
head\_left: $<$image$>$
cam\_wrist\_left: $<$image$>$
cam\_wrist\_right: $<$image$>$
Task: $<$instruction$>$
State: $<$$\vert$robot\_state\_implicit\_stats$\vert$$>$$<$$\vert$robot\_state$\vert$$>$
Control Mode: Arms: EE, LOW BODY: JOINT
}

\begin{itemize}
\item \texttt{arm\_type == "dual\_with\_legs"} 才加 \texttt{, LOW BODY: JOINT}
\item 图像按 \texttt{image\_roles} 顺序，\textbf{role 标签紧跟在自己的图前面}
\item 注释明确说这个格式是为了*\emph{对齐离线预训练数据构建脚本 \texttt{build\_vlm\_dataset.py}*} 的输入分布——
即 VLM 底座在预训练阶段就见过这套 prompt，这里必须复现，否则分布外
\end{itemize}

\textbf{(3) 多模态占位符检测的坑}（\texttt{QWen3.py:148-154}）
必须保留 \texttt{input\_ids} 而不能只传 \texttt{inputs\_embeds}。因为 Qwen3VL 的 \texttt{get\_placeholder\_mask} 用
\textbf{token id 比较}定位图像 token；若丢掉 \texttt{input\_ids}，它会退化到"比较 embedding 是否相等"的分支，
一旦别的 special token 的 embedding 与 image-pad embedding 碰撞就会误计数，报
\texttt{Image features and image tokens do not match}。保留 \texttt{input\_ids} 也让 M-RoPE 能正常计算。

\subsection*{3.3 robot\_state 注入：\texttt{projector/robot\_state.py} + hook}

有两种把机器人状态喂给模型的通路，二者\textbf{互斥}：

$\vert$ $\vert$ 通路 A：DiT 侧 state token $\vert$ 通路 B：VLM 侧 projector $\vert$
$\vert$---$\vert$---$\vert$---$\vert$
$\vert$ 开关 $\vert$ \texttt{robot\_state\_projector.enabled = false}（默认） $\vert$ \texttt{enabled = true} $\vert$
$\vert$ 实现 $\vert$ \texttt{state\_encoder} MLP，把 \texttt{state⊕mask}(120) → 768，作为\textbf{第 0 个 token} 拼在动作 token 前 $\vert$ \texttt{RobotStateProjector} 把 120 维投影到 VLM hidden(2048)，\textbf{散射到 \texttt{$<$$\vert$robot\_state$\vert$$>$} token 位置} $\vert$
$\vert$ 代码 $\vert$ \texttt{MMDiT\_ActionHeader.py:139-147, 266-270} $\vert$ \texttt{QWen3.py:95-196} $\vert$
$\vert$ 副作用 $\vert$ 无 $\vert$ \texttt{QwenMMDiT.py:199-200} 会强制 \texttt{state\_dim = 0}，关掉通路 A $\vert$

通路 B 的实现很巧妙（\texttt{QWen3.py:156-196}）：它\textbf{不给 embedding 层加 forward hook}，
而不是预先算好 \texttt{inputs\_embeds} 再置空 \texttt{input\_ids}。这样做的好处是
① 保住 token-id 占位符检测，② 保住 M-RoPE，③ \texttt{masked\_scatter} 可微所以 projector 能端到端训练。
hook 里还做了严格校验：\texttt{$<$$\vert$robot\_state$\vert$$>$} 出现次数必须等于 batch size，否则直接报错。

\texttt{RobotStateProjector} 本身就是一个 2 层 MLP：\texttt{Linear(120→hidden) → SiLU → Linear(hidden→VLM\_hidden)}。

\subsection*{3.4 动作头：\texttt{MMDiT\_ActionHeader.py}}

\subsubsection*{3.4.1 配置的"三次覆盖"}

\texttt{\_\_init\_\_} (\texttt{MMDiT\_ActionHeader.py:106-121}) 里有个容易看晕的层叠：

\texttt{python
config = full\_config.framework.action\_model       \# YAML 的值
action\_model\_cfg = DiTConfig["DiT-B"]             \# 硬编码 \{"input\_embedding\_dim":768,
                                                  \#           "attention\_head\_dim":64,
                                                  \#           "num\_attention\_heads":12\}
config.update(action\_model\_cfg)                   \# ① 覆盖 action\_model 节点
config.diffusion\_model\_cfg.update(action\_model\_cfg)  \# ② 覆盖 diffusion\_model\_cfg
diffusion\_model\_cfg = \{**action\_model\_cfg, **diffusion\_model\_cfg\}  \# ③ YAML 优先
self.model = DiT(**diffusion\_model\_cfg)
}

最终生效的 DiT 形状（用仓库自带 YAML）：

$\vert$ 参数 $\vert$ 值 $\vert$ 来源 $\vert$
$\vert$---$\vert$---$\vert$---$\vert$
$\vert$ \texttt{num\_attention\_heads} $\vert$ 12 $\vert$ DiTConfig（YAML 没写） $\vert$
$\vert$ \texttt{attention\_head\_dim} $\vert$ 64 $\vert$ DiTConfig $\vert$
$\vert$ \texttt{inner\_dim} $\vert$ \textbf{768} = 12×64 $\vert$ 计算得出 $\vert$
$\vert$ \texttt{cross\_attention\_dim} $\vert$ \textbf{2048} $\vert$ 运行时从 VLM hidden 对齐（\texttt{QwenMMDiT.py:204-206}） $\vert$
$\vert$ \texttt{num\_layers} $\vert$ 16 $\vert$ YAML $\vert$
$\vert$ \texttt{output\_dim} $\vert$ \textbf{1024} $\vert$ YAML $\vert$

*\emph{注意 \texttt{input\_embedding\_dim=768} 但 \texttt{output\_dim=1024}*}，两者不等是刻意的：
动作 token 在 768 维空间里做联合注意力（和 12 头×64 匹配），最后由 \texttt{action\_decoder}
\texttt{MLP(1024→1024→54)} 投到动作维度。

\subsubsection*{3.4.2 网络组件}

\texttt{
                 ┌─ (可选) state\_encoder: MLP(120 → 1024 → 768) ──┐  1 个 token
                 │                                                │
noisy\_action ──▶ ActionEncoder ──▶ +pos\_embed ──▶ FiLM\_pre ──────┴──▶ hidden\_states (B, 1+30, 768)
(B,30,108)       Linear(108→768)   Embedding(1024,768)                      │
                 + 正弦时间编码      (按位置 0..29)                          │
                 swish/Linear                                               ▼
                                                              ┌──────────────────────────┐
VLM last\_hidden ─────────────────────────────────────────────▶│  DiT (16 × QwenImageBlock)
(B, L, 2048)     RMSNorm → Linear(2048→768)                   │  双流联合注意力 + RoPE   │
                                                              └───────────┬──────────────┘
                                                                          ▼
                                                       FiLM\_dec ──▶ action\_decoder
                                                       MLP(1024→1024→54) ──▶ (B,30,54)
}

\begin{itemize}
\item \texttt{ActionEncoder}（\texttt{MMDiT\_ActionHeader.py:27-52}）：把 \texttt{[noisy\_action ‖ mask]}(108) 线性到 768，
与\textbf{离散化时间步}的正弦编码拼接，过 swish + Linear。注意时间用的是
\texttt{t\_discretized = (t * 1000).long()}（整数桶）而非连续 t。
\item \texttt{state\_encoder} 用 \texttt{MLP(120→1024→768)}：输入是 \textbf{60 维 state ⊕ 60 维 mask = 120 维}，
与文档 §4.5 一致。
\item \texttt{embodiment\_embedding} + \texttt{FiLM}（\texttt{:163-214}）：按 \texttt{robot\_type} 查表得到 embedding，
用 \texttt{Linear} 投成 \texttt{(scale, shift)}，以 \texttt{x*(1+scale)+shift} 施加在\textbf{两处}：
action token 进 DiT 前（768 维）、DiT 输出进 decoder 前（1024 维）。
两个 FiLM 都是 \textbf{zero-init}（AdaLN-Zero），所以训练起步时等价于"没有本体条件"。
为什么需要它？见 §6.3。
\end{itemize}

\subsubsection*{3.4.3 流匹配（flow matching）的完整数学}

\textbf{训练}（\texttt{forward}, \texttt{MMDiT\_ActionHeader.py:230-291}）：

\texttt{python
t          = noise\_s * (1 - Beta(1.5, 1.0).sample())      \# t ∈ (0, 0.999]
noisy      = (1 - t) * noise + t * action                 \# 线性插值路径
velocity   = action - noise                               \# 回归目标
pred       = decoder(DiT(action\_encoder([noisy ‖ mask]), VLM\_hidden, t))
sq\_err     = (pred - velocity) ** 2
loss       = (sq\_err * mask).sum() / mask.sum()            \# 掩码加权 MSE
}

\texttt{Beta(1.5,1.0)} 的密度 ∝ x\textasciicircum{}0.5，质量偏向 x→1，所以 \texttt{t} 偏向 0，即\textbf{采样偏向高噪声端}。
\texttt{mask\_part} 会在算 \texttt{noisy} 后被乘回去，注释解释这是为了修掉"训练时编码器看到非零无效维、
推理时看到零"的输入不一致。

\textbf{推理}（\texttt{predict\_action}, \texttt{:293-350}）：从纯噪声出发，4 步 Euler 积分

\texttt{python
dt = 1/num\_inference\_timesteps            \# 0.25
for step in range(4):
    v = model(x\_t, cond, t=step/4)
    x\_t = (x\_t + dt * v) * mask           \# 无效维每步清零
}

\textbf{\texttt{repeated\_diffusion\_steps = 8}**（\texttt{QwenMMDiT.py:289-305}）：把同一批样本重复 8 次，
让每个样本贡献 8 个}独立采样**的 (t, noise) 组合。这样单向 VLM 前向只跑 1 次，
却能拿到 8 倍流匹配监督信号——用显存换梯度质量的标准做法。

\subsection*{3.5 DiT 主干：\texttt{DiT\_modules/mmdit.py}}

这是从 Qwen-Image 移植的双流（double-stream）MMDiT。

\subsubsection*{3.5.1 双流联合注意力（\texttt{QwenDoubleStreamAttnProcessor2\_0}, \texttt{:169-305}）}

两个"流"：\textbf{action 流}（\texttt{to\_q/k/v}）与 \textbf{text 流}（\texttt{add\_q\_proj/add\_k\_proj/add\_v\_proj}）。
每层做的事情：

\texttt{
1. 分别算两流的 QKV，各自过 QK-norm              (:218-226)
2. 分别施加 RoPE                                  (:229-234)
3. 拼接: joint = [action, text]   ← action 在前!   (:248-250)
4. 一次联合注意力（flash / flash\_varlen / native） (:271-285)
5. 切回两流，各自输出投影 + 残差                    (:295-303)
}

\subsubsection*{3.5.2 两个"必须 action-first"的理由（注释写得很清楚，\texttt{mmdit.py:236-247} 与 \texttt{:545-556}）}

\begin{enumerate}
\item \textbf{FLASH\_VARLEN 的 key packer 会 \texttt{key[b, :valid\_len]} 从前端切片}，
它默认"有效位置是从 0 开始的一段连续前缀"。文本流里被 mask 掉的
assistant turn / 右 padding 会造成中间空洞；把永远有效的 action 放最前面，
就能保证 True 区间是 0-anchored 的连续前缀。顺序反了会\textbf{静默地注意错 key}。
\item \texttt{QwenActionTextRotaryEmbed}（\texttt{:46-85}）本身也按 action-first 分配位置：
action 占 \texttt{0..S\_act-1}，text 从 \texttt{S\_act} 续接。注释解释了为什么这样锚定：
\textbf{action 块长度固定（= action\_horizon），text 长度在训练/推理间变化}，
把 action 锚在 0 可以让动作流的旋转位置在训练和推理时完全一致。
\end{enumerate}

\subsubsection*{3.5.3 工程细节：为导出 ONNX 而写的代码}

代码里有好几处注释指明某个写法是为了让 torch dynamo / ONNX exporter 正确 trace：

\begin{itemize}
\item \texttt{pos = torch.arange(total, dtype=torch.float32)}（而非默认 int64）——
旧 ONNX exporter 会误把 \texttt{torch.outer} 的类型提升 trace 成 int64 Sin/Cos 节点（\texttt{:76-79}）
\item 用 \texttt{reshape(*x.shape[:-1], -1, 2)} 而非 \texttt{.flatten(3)}——
\texttt{.flatten(3)} 被 lower 成把 trace 时 batch 维写死的 Reshape（\texttt{:147-150}）
\item 用实数版 RoPE（\texttt{use\_real=True}）而非 \texttt{view\_as\_complex}——复数张量在 ONNX/TensorRT 里没有表示（\texttt{:62-65}）
\item \texttt{\_attention\_backend} 可切到 \texttt{AttentionBackendName.NATIVE}（纯 SDPA）以便导出，
因为 flash-attn kernel 没有 ONNX 等价物（\texttt{:252-258}）
\end{itemize}

\textbf{这告诉你：这个模型是要上真机/上推理引擎的}，不只是一篇论文的参考实现。

\subsubsection*{3.5.4 AdaLN 调制}

\texttt{QwenImageTransformerBlock}（\texttt{:307-444}）每层有两个调制：
\texttt{action\_mod} / \texttt{txt\_mod} 各把 \texttt{temb} 投成 \texttt{6*dim}，切成 2 份 \texttt{3*dim}（norm1 用、norm2 用），
每份再切成 \texttt{(shift, scale, gate)}，形式为 \texttt{x*(1+scale)+shift} 配 \texttt{x + gate*sublayer(x)}。
\texttt{\_modulate} 支持 \texttt{[B,D]}（广播到整条序列）和 \texttt{[B,L,D]}（RTC 逐步调制）两种形状——
\textbf{RTC（real-time chunking）是预留但当前未启用的路径}，标准调度下 \texttt{temb} 是 \texttt{[B,D]}。

\texttt{DiT.forward} 的最终 AdaLN 用的是 \texttt{temb} 而非 \texttt{txt\_temb}（\texttt{:585}），因为要调制的是 action 流。

\subsection*{3.6 组装者：\texttt{framework/VLM4A/QwenMMDiT.py}}

\subsubsection*{\texttt{forward}（训练，\texttt{:233-318}）}

```python
examples = \emph{adapt}multi\emph{source}batch(batch)     \# collate 的 dict → List[dict]（含 PIL 图）
qwen\emph{inputs = build}qwenvl\emph{inputs(images, instructions, solutions=None, image}roles, arm\emph{types)
backbone}attention\emph{mask = qwen}inputs.get("attention\emph{mask")
assistant}mask = qwen\emph{inputs.pop("assistant}mask", None)

last\emph{hidden = qwen}vl\emph{interface(**qwen}inputs, robot\emph{states=...).hidden}states[-1]   \# (B, L, 2048)

\chapter*{构造 DiT 的 key 掩码：文本按 attention\_mask，assistant turn 要被排除}

keys\emph{attn}mask = backbone\emph{attention}mask \& \textasciitilde{}assistant\_mask

actions\emph{target = actions[:, -action}horizon:, :]                     \# (B, 30, 54)
actions\emph{target = cat([actions}target, mask.expand\_as(...)], dim=2)   \# (B, 30, 108)

\chapter*{8 倍重复}

last\emph{hidden}r  = last\emph{hidden.repeat(8, 1, 1)
actions}r      = actions\emph{target.repeat(8, 1, 1)
state}r        = state.repeat(8, 1, 1)
body\emph{ids}r     = body\_ids.repeat(8)

flow\emph{action}loss = action\emph{model(last}hidden\emph{r, actions}r, state\emph{r, keys}attn\emph{mask}r, body\emph{ids}r)
return \{"flow\emph{action}loss": loss, "action\_loss": loss\}
```

\subsubsection*{\texttt{\_build\_state\_tensor}（\texttt{:69-106}）—— 训练期状态增强}

$\vert$ 概率 $\vert$ 行为 $\vert$
$\vert$---$\vert$---$\vert$
$\vert$ 10\% $\vert$ state 和 state\_mask \textbf{全部清零} $\vert$
$\vert$ 30\% $\vert$ 给有效维加 \texttt{N(0, 0.1)} 噪声 $\vert$
$\vert$ 60\% $\vert$ 保持原样 $\vert$

这是一种廉价而有效的鲁棒性正则：让模型学会在"状态缺失/有噪"时依然能靠视觉完成任务。
\textbf{只在 \texttt{training=True} 时启用}。

\subsubsection*{\texttt{predict\_action}（推理，\texttt{:320-378}）}

与 forward 的差异：
- \texttt{add\_generation\_prompt=False}（不留 \texttt{$<$$\vert$im\_start$\vert$$>$assistant\textbackslash{}n}）
- 用 \texttt{to\_pil\_preserve} 而非 \texttt{TF.to\_pil\_image}
- 若配置了 \texttt{obs\_image\_size} 会 \texttt{resize\_images} 对齐训练分辨率
- 不构造 action\_mask 的重复，直接走 4 步 Euler

\hrule
\section*{4. 训练层}

\subsection*{4.1 启动流程（\texttt{train\_unifolm\_wla.py:478-508} → \texttt{main():441-475}）}

```python
cfg = OmegaConf.load(args.config\emph{yaml)
cfg = OmegaConf.merge(cfg, OmegaConf.from}dotlist(normalize\emph{dotlist}args(clipargs)))
cfg = apply\emph{config}compat(cfg)          \# 老配置 → version\_id "0.21" schema
main(cfg)

\chapter*{main():}

cfg = wrap\emph{config(cfg)                  \# 包成 AccessTrackedConfig
accelerator = Accelerator(gradient}accumulation\emph{steps=..., deepspeed}plugin=DeepSpeedPlugin())
output\emph{dir = setup}directories(cfg)
vla = build\emph{framework(cfg)              \# 建模型（预训练底座在这里加载）
vla}train\emph{dataloader = prepare}data(...)
optimizer, lr\emph{scheduler = setup}optimizer\emph{and}scheduler(vla, cfg)
trainer = VLATrainer(...); trainer.prepare\_training(); trainer.train()
```

\texttt{apply\_config\_compat()}（\texttt{share\_tools.py:418-524}）是一个\textbf{幂等的配置规范化层}，
负责把用户 YAML 的简写补全成下游代码期待的形状：

$\vert$ 规则 $\vert$ 说明 $\vert$
$\vert$---$\vert$---$\vert$
$\vert$ \texttt{action\_horizon} ↔ \texttt{future\_action\_window\_size} $\vert$ 前者是规范名，自动补 \texttt{future = horizon - 1}；不一致时 warning 并以 \texttt{action\_horizon} 为准 $\vert$
$\vert$ \texttt{diffusion\_model\_cfg.output\_dim} $\vert$ 缺省从 \texttt{action\_model.hidden\_size} 补 $\vert$
$\vert$ \texttt{diffusion\_model\_cfg.cross\_attention\_dim} $\vert$ 缺省从 \texttt{framework.qwenvl.vl\_hidden\_dim} 补 $\vert$
$\vert$ \texttt{action\_model.action\_hidden\_dim} $\vert$ 缺省从 \texttt{hidden\_size} 补 $\vert$
$\vert$ \texttt{past\_action\_window\_size} $\vert$ 缺省补 0（所有已发布框架都是 0，所以从用户 YAML 里删掉了） $\vert$
$\vert$ \texttt{version\_id} $\vert$ 盖戳 \texttt{"0.21"} $\vert$

\subsection*{4.2 \texttt{AccessTrackedConfig}：把"用到的配置"写出来（\texttt{config\_tracker.py}）}

这个模块的动机很实际：\textbf{训练配置有几百个字段，但一次运行真正影响结果的只有其中一部分}。
它用 \texttt{\_\_getattr\_\_} / \texttt{\_\_getitem\_\_} / \texttt{get()} 记录每一次访问，最后
\texttt{export\_accessed\_config()} 只把访问到的\textbf{叶子路径}写进 \texttt{config.yaml}。
好处是产出的 checkpoint 目录自带一份"最小可复现配置"，\texttt{from\_pretrained} 时能直接重建模型。

三个实现细节：

\begin{enumerate}
\item \texttt{\_children} 缓存子 config 的 wrapper；\texttt{\_\_setattr\_\_} 时必须\textbf{失效缓存}，
否则会拿到指向旧节点的 stale wrapper。
\item \texttt{merge\_framework\_config}（\texttt{share\_tools.py:230-248}）在替换 \texttt{framework} 节点后，
会*\emph{保留并恢复 \texttt{\_local\_accessed}*}，否则 \texttt{build\_framework} 里读过的 \texttt{name} 字段
会从 \texttt{config.yaml} 里神秘消失。
\item 文件末尾（\texttt{config\_tracker.py:486-549}）\textbf{monkey patch 了 5 个 OmegaConf 全局函数}
（\texttt{to\_container} / \texttt{save} / \texttt{to\_yaml} / \texttt{is\_config} / \texttt{merge}）让它们能接受
\texttt{AccessTrackedConfig}。这是个有争议的做法——它会把影响扩散到所有 import 了 omegaconf 的代码。
\end{enumerate}

\subsection*{4.3 \texttt{VLATrainer.prepare\_training()}（\texttt{:120-159}）}

顺序很重要：

\begin{enumerate}
\item \texttt{set\_seed(cfg.seed + rank)} —— 每卡不同种子
\item \texttt{\_save\_initial\_configs()} —— \textbf{先落盘 config.full.yaml 和 config.yaml}。
注释说明动机：万一后面 checkpoint 加载 / DeepSpeed 初始化 / dataloader 构建崩溃，
产出的 run 目录依然是可自省、可 \texttt{from\_pretrained} 的
\item \texttt{\_init\_checkpointing()} —— \texttt{is\_resume} 时找最新 checkpoint，否则加载 \texttt{pretrained\_checkpoint}
\item \texttt{\_adjust\_lr\_scheduler\_for\_resume()} —— 用 \texttt{for \_ in range(completed\_steps): lr\_scheduler.step()} 补步
\item \texttt{freeze\_backbones()} 按 \texttt{freeze\_modules}（逗号分隔的模块路径）冻结
\item \texttt{print\_trainable\_parameters()}
\item \texttt{accelerator.prepare(model, optimizer[, dataloader])}
\item \texttt{\_init\_wandb()}（\texttt{mode="offline"}）
\end{enumerate}

\subsection*{4.4 训练主循环（\texttt{:308-352}）}

\texttt{python
while completed\_steps $<$ max\_train\_steps:
    batch = self.\_get\_next\_batch()        \# StopIteration 时自动重置 iterator + set\_epoch
    step\_metrics = self.\_train\_step(batch)
    if accelerator.sync\_gradients:
        completed\_steps += 1              \# 只在梯度同步时计步
    if completed\_steps \% eval\_interval == 0:  step\_metrics = eval\_action\_model(step\_metrics)
    \_log\_metrics(step\_metrics)
    if completed\_steps \% save\_interval == 0 and completed\_steps $>$ 0: \_save\_checkpoint()
}

\texttt{\_train\_step}（\texttt{:385-416}）：

\texttt{python
with accelerator.accumulate(model):
    optimizer.zero\_grad()
    with torch.autocast("cuda", dtype=torch.bfloat16):
        output\_dict = model.forward(batch)
        total\_loss = output\_dict["action\_loss"]
    accelerator.backward(total\_loss)
    if gradient\_clipping: accelerator.clip\_grad\_norm\_(...)
    optimizer.step()
    if accelerator.sync\_gradients:        \# ← 关键守卫
        lr\_scheduler.step()
}

\texttt{sync\_gradients} 守卫的注释解释得很清楚：不加它，scheduler 会以
\texttt{gradient\_accumulation\_steps} 倍速前进，导致 warmup 提前结束、cosine 在
\texttt{max\_train\_steps} 之前就衰减到 \texttt{min\_lr}。

\subsection*{4.5 优化器与分层学习率（\texttt{trainer\_tools.py:88-144}）}

\texttt{build\_param\_lr\_groups} 读 \texttt{cfg.trainer.learning\_rate} 这个 dict：

\texttt{yaml
learning\_rate:
  base: 1.0e-04                  \# 兜底
  qwen\_vl\_interface: 1.0e-05     \# VLM 用 1/10 学习率
  action\_model: 1.0e-04          \# 动作头全速
}

对每个非 \texttt{base} 的 key，用 \texttt{getattr} 逐级解析模块路径，取其参数建一个 param group；
剩余参数归入 \texttt{base} 组。\texttt{freeze\_modules} 里的参数会被排除。
这实现了「VLM 小步微调、动作头从零训练」的经典配比。

\subsection*{4.6 日志与 checkpoint}

\begin{itemize}
\item \texttt{initialize\_overwatch()}（\texttt{overwatch.py:146}）返回 \texttt{DistributedOverwatch}（主进程 INFO、
其他进程 ERROR）或 \texttt{PureOverwatch}（单进程 / spawn worker）
\item checkpoint 命名 \texttt{steps\_\{N\}\_pytorch\_model.pt} 或 \texttt{steps\_\{N\}\_model.safetensors}，
\texttt{\_get\_latest\_checkpoint} 用正则严格匹配这个名字（\texttt{trainer\_tools.py:520-522}）
\item 训练结束额外存一份 \texttt{final\_model/}
\item 同时追加写 \texttt{summary.jsonl}
\end{itemize}

\subsection*{4.7 训练 YAML 与 DeepSpeed 配置}

\texttt{config/training/unifolm\_wla\_pretrain\_multisource.yaml} 关键值：
\texttt{action\_horizon: 30}、\texttt{action\_dim: 54}、\texttt{state\_dim: 60}、\texttt{repeated\_diffusion\_steps: 8}、
\texttt{per\_device\_batch\_size: 8}、\texttt{gradient\_accumulation\_steps: 4}、\texttt{max\_train\_steps: 200000}、
\texttt{num\_warmup\_steps: 2000}、\texttt{save\_format: safetensors}、\texttt{freeze\_modules: ''}（全量训练）。

\texttt{config/deepseeds/} 下有 6 个配置，实际是两类：

\begin{itemize}
\item \texttt{ds\_config.yaml} / \texttt{ds\_config\_debug\_cpu\_offload.yaml} / \texttt{zero3.yaml} —— \textbf{纯 DeepSpeed json}
\item \texttt{deepspeed\_zero2.yaml} / \texttt{deepspeed\_zero3.yaml} / \texttt{zero2.yaml} —— \textbf{Accelerate 的启动配置}
（\texttt{distributed\_type: DEEPSPEED} + 指向上面的 json）
\end{itemize}

生产用 \texttt{ds\_config.yaml}：\textbf{ZeRO stage 2}（不是 3）、\texttt{bf16: true}、\texttt{allgather/reduce\_bucket\_size: 5e8}、
\texttt{overlap\_comm} + \texttt{contiguous\_gradients}。stage 2 而非 3 的选择是合理的——
VLM 4B 参数在 8 卡上 ZeRO-2 足够，且省掉 stage-3 的反复 all-gather 开销。

\hrule
\section*{5. 端到端 Trace（手算维度）}

以仓库默认配置（\texttt{action\_horizon=30}、DiT-B、Qwen3-VL-4B、batch=1）为例：

$\vert$ \# $\vert$ 位置 $\vert$ 形状 $\vert$ 说明 $\vert$
$\vert$---$\vert$---$\vert$---$\vert$---$\vert$
$\vert$ 1 $\vert$ \texttt{LeRobotDataset.\_\_getitem\_\_} $\vert$ action \texttt{(30, D\_raw)}, state \texttt{(D\_raw,)} $\vert$ D\_raw 因数据集而异 $\vert$
$\vert$ 2 $\vert$ \texttt{compute\_relative\_actions} $\vert$ \texttt{(30, 6)} $\vert$ 左/右臂各一份 $\vert$
$\vert$ 3 $\vert$ \texttt{map\_action\_chunk} $\vert$ \texttt{(30, 54)} + \texttt{(54,)} bool $\vert$ 统一槽位 $\vert$
$\vert$ 4 $\vert$ \texttt{map\_state} $\vert$ \texttt{(60,)} + \texttt{(60,)} bool $\vert$ rot6d 布局 $\vert$
$\vert$ 5 $\vert$ \texttt{collate\_fn} $\vert$ action \texttt{(B,30,54)}, state \texttt{(B,60)}, images \texttt{\{role:(B,3,336,448)\}} $\vert$ $\vert$
$\vert$ 6 $\vert$ \texttt{\_adapt\_multi\_source\_batch} $\vert$ \texttt{List[B]} of \texttt{\{image:[PIL], lang, action, state, ...\}} $\vert$ 回到样本级 $\vert$
$\vert$ 7 $\vert$ \texttt{build\_qwenvl\_inputs} $\vert$ \texttt{input\_ids (B,L)}, \texttt{pixel\_values}, \texttt{attention\_mask} $\vert$ L ≈ 3×\textasciitilde{}100(图) + \textasciitilde{}40(文本) $\vert$
$\vert$ 8 $\vert$ \texttt{qwen\_vl\_interface(...)} $\vert$ \texttt{hidden\_states[-1]} = \texttt{(B, L, 2048)} $\vert$ 输出 hidden 而非 logits $\vert$
$\vert$ 9 $\vert$ \texttt{\_build\_state\_tensor} $\vert$ \texttt{(B, 1, 120)} $\vert$ 60 state ⊕ 60 mask $\vert$
$\vert$ 10 $\vert$ \texttt{actions\_target} + mask 拼接 $\vert$ \texttt{(B, 30, 108)} $\vert$ 54 动作 ⊕ 54 mask $\vert$
$\vert$ 11 $\vert$ ×8 repeat $\vert$ \texttt{(8B, 30, 108)}, \texttt{(8B, L, 2048)} $\vert$ $\vert$
$\vert$ 12 $\vert$ \texttt{ActionEncoder} $\vert$ \texttt{(8B, 30, 768)} $\vert$ 108 → 768 $\vert$
$\vert$ 13 $\vert$ \texttt{state\_encoder} $\vert$ \texttt{(8B, 1, 768)} $\vert$ 120 → 1024 → 768 $\vert$
$\vert$ 14 $\vert$ \texttt{hidden\_states} $\vert$ \texttt{(8B, 31, 768)} $\vert$ 1 state token + 30 action tokens $\vert$
$\vert$ 15 $\vert$ \texttt{DiT} 输入投影 $\vert$ text \texttt{(8B, L, 2048)} → RMSNorm → Linear → \texttt{(8B, L, 768)} $\vert$ $\vert$
$\vert$ 16 $\vert$ 16× 双流联合注意力 $\vert$ action 流 \texttt{(8B,31,768)}，text 流 \texttt{(8B,L,768)} $\vert$ joint 长度 \texttt{31+L} $\vert$
$\vert$ 17 $\vert$ \texttt{proj\_out\_2} $\vert$ \texttt{(8B, 31, 1024)} $\vert$ 768 → 1024 $\vert$
$\vert$ 18 $\vert$ \texttt{action\_decoder} $\vert$ \texttt{(8B, 31, 54)} $\vert$ 1024 → 1024 → 54 $\vert$
$\vert$ 19 $\vert$ \texttt{pred\_actions = pred[:, -30:]} $\vert$ \texttt{(8B, 30, 54)} $\vert$ \textbf{取尾部 30 个 token} $\vert$
$\vert$ 20 $\vert$ 掩码加权 MSE vs \texttt{velocity} $\vert$ scalar $\vert$ $\vert$

\textbf{第 19 步值得注意}：因为第 14 步把 state token 放在\textbf{最前面}，所以"最后 30 个"正好是 action 部分。
这是 \texttt{pred[:, -H:]} 而非 \texttt{pred[:, :H]} 的原因。

\hrule
\section*{6. 关键设计决策：为什么这么做}

\subsection*{6.1 为什么 VLM 只当编码器，不做自回归生成？}

\texttt{QwenMMDiT.forward} 调用 \texttt{build\_qwenvl\_inputs(..., solutions=None)}，
于是 \texttt{QWen3.py:316} 的 \texttt{if solutions is not None} 分支不进入 →
*\emph{不构造 \texttt{labels}，不算 LM-CE loss，没有 \texttt{assistant\_mask}*}。
VLM 的梯度完全来自 DiT 的 cross-attention。

好处：一次前向拿到全部视觉/语言特征，不需要 teacher-forcing 自回归解码；
省掉 LM head 和 tokenizer 的监督开销；动作输出天然是连续、多模态的（flow matching 擅长这个）。

代价：失去"让 VLM 同时学会说人话/做 CoT"的能力。代码里保留了 \texttt{forward\_vlm()} 和
\texttt{compute\_loss(tag=...)} 这套多任务接口（\texttt{base\_framework.py:137-203}）作为将来扩展点，
但\textbf{当前 trainer 从不调用它们}（直接 \texttt{model.forward(batch)}）。

\subsection*{6.2 为什么用 MMDiT 双流而不是简单的 cross-attention？}

标准 DiT 的 cross-attention 里文本流是被动的（只提供 K/V，自身不更新）。
MMDiT 双流让\textbf{文本/视觉流自己也过 MLP 和残差}，在每一层都被动作流的信息"回流"影响。
\texttt{DiT.forward} 虽然最终丢弃了 text 流（只返回 action 流），但 text 流在中间层是被持续更新的。
这对"动作 token 数量少（31）、条件 token 多（L）"的场景有意义。

\subsection*{6.3 为什么需要 embodiment embedding？（\texttt{MMDiT\_ActionHeader.py:82-97} 的注释）}

注释直白地说明了问题：

\begin{quote}
Single-arm sources (bridge/fractal/libero/...) all map into the SHARED left action slot,
so without an embodiment signal the shared decoder averages their conflicting distributions
and \textbf{the left arm collapses to \textasciitilde{}0}.
\end{quote}

即：多个单臂机器人共用左臂槽位，各自的动作分布互相冲突，共享解码器会学成"平均值≈0"。
解法是给 DiT 一个 \texttt{robot\_type → embedding} 的条件信号，让共享的 encoder/decoder 保持本体无关，
由 DiT 去消歧。\texttt{DEFAULT\_EMBODIMENT\_TYPES} 的顺序是\textbf{序列化契约}，必须 append-only，
且 index 0 保留给 "unknown"（\texttt{build\_embodiment\_map} 返回 \texttt{i+1}）。

\subsection*{6.4 为什么 DiT 的 cross\emph{attention}dim 要在运行时对齐？}

\texttt{QwenMMDiT.py:204-206} 用 \texttt{self.qwen\_vl\_interface.model.config.hidden\_size} 覆盖 YAML 值。
因为 VLM 底座可能被换（Qwen3-VL-2B/4B/8B hidden 分别是 2048/2048/4096 级别），
硬编码会让换底座时静默出错。运行时读实际值最安全。

\subsection*{6.5 为什么用 \texttt{spawn} 而不是 \texttt{fork}？}

\texttt{dataloader.py:203-205} 注释：fork 出来的 worker 会继承 trainer 已经初始化的 CUDA context。
数据通路是 CPU-only，继承 CUDA context 纯属浪费且有 fork+CUDA 的未定义行为风险。

\hrule
\section*{7. 坑与风险清单}

以下是我通读后认为\textbf{值得作者确认}的点。按严重程度排序。

\subsection*{7.1 🔴 \texttt{config/training/*.yaml} 的 \texttt{framework.name} 是 \texttt{QwenGR00T}，但只注册了 \texttt{QwenMMDiT}}

\texttt{unifolm\_wla\_pretrain\_multisource.yaml:9} → \texttt{QwenGR00T}。
\texttt{build\_framework} 会抛 \texttt{NotImplementedError: Framework 'QwenGR00T' is not implemented.
Available frameworks: ['QwenMMDiT']}。\textbf{仓库自带的训练配置开箱即跑不通}，需要手改。

同时 \texttt{configs/unitree.yaml:121} 的 WBT 数据源 \texttt{data\_path} 也写成了 \texttt{UnifoLM\_G1\_Dex1\_Dataset/}，
而 \texttt{docs/train\_action\_expert.md:41,50} 明确要求 WBT 用 \texttt{UnifoLM\_WBT\_Dataset}。
看起来是复制粘贴残留，会让两个数据源指向同一份数据。

\subsection*{7.2 🟠 训练/推理的 DiT key 集合不一致（尾部差一个 generation prompt）}

推导链条：

\begin{enumerate}
\item \texttt{QwenMMDiT.forward} 传 \texttt{solutions=None} → \texttt{add\_generation\_prompt = (solutions is None) = True}
（\texttt{QWen3.py:300-301}）
\item 因此\textbf{训练}时 prompt 尾部是 \texttt{...Control Mode: ...$<$$\vert$im\_end$\vert$$>$\textbackslash{}n$<$$\vert$im\_start$\vert$$>$assistant\textbackslash{}n}
\item \texttt{assistant\_mask} 是 \texttt{None}（\texttt{:316} 的 \texttt{if solutions is not None} 不成立），
所以 \texttt{QwenMMDiT.py:272-273} 的 \texttt{\textasciitilde{}assistant\_mask} 掩码逻辑\textbf{从未生效}
\item \texttt{predict\_action} 传 \texttt{add\_generation\_prompt=False}（\texttt{QwenMMDiT.py:341}）
→ \textbf{推理}时 key 集合停在 \texttt{...$<$$\vert$im\_end$\vert$$>$}
\end{enumerate}

结果是训练时 DiT 多看了末尾 3 个 token。\texttt{QWen3.py:293-299} 的注释其实假设了
"训练时有 solutions，assistant turn 被 mask 掉，所以推理去掉 header 正好对齐"——
但那套假设在 QwenMMDiT 这条路径上不成立。

影响量级不大（3 个 key 位置），但这是\textbf{训练/推理分布不一致}，属于应该修的问题。
修法有两种：训练时也传 \texttt{add\_generation\_prompt=False}，或让 \texttt{forward} 显式传空 solutions
并调整分支。

\subsection*{7.3 🟠 \texttt{action\_horizon} 与实际 chunk 长度的静默耦合}

\texttt{actions\_target = actions[:, -self.action\_horizon:, :]}（\texttt{QwenMMDiT.py:277}）\textbf{取尾部}。
而 chunk 长度等于 \texttt{source\_fps}（=30），\texttt{chunk\_size} 配置项从未被用于切片。
当前 YAML 里 \texttt{action\_horizon=30} 恰好等于 \texttt{source\_fps=30}，所以没暴露问题。
但如果有人把 \texttt{action\_horizon} 调成 8，会静默地只拿 \texttt{t=22..29} 这 8 步（最后一秒的后 1/4），
而不是直觉上的"前 8 步"。\textbf{建议加断言} \texttt{chunk\_len $>$= action\_horizon} 并明确语义。

\subsection*{7.4 🟠 \texttt{target\_fps} / \texttt{source\_fps} 不一致时不会重采样}

\texttt{fps\_utils.py} 里的 \texttt{resample\_action\_chunk} 和 \texttt{build\_bspline\_resample\_matrix}
\textbf{没有任何调用点}。\texttt{\_build\_delta\_timestamps} 只用 \texttt{source\_fps} 生成时间戳，
\texttt{target\_fps} 仅出现在一条 log 里（\texttt{single\_source\_dataset.py:118}）。
所以混入一个 10 FPS 的数据源时，它的 chunk 是 10 步、时间跨度 1 秒，
而 30 FPS 的是 30 步、1 秒 —— 两者在 \texttt{action\_horizon=30} 的约束下行为完全不同
（10 步的会被 LeRobot 的 delta 机制 pad/复制到 30 个查询点？需要实测确认）。

\subsection*{7.5 🟡 \texttt{eval\_action\_model} 会吃掉一个训练 batch}

\texttt{train()} 里每 \texttt{eval\_interval} 步调用 \texttt{eval\_action\_model()}，而它内部
\texttt{examples = self.\_get\_next\_batch()}（\texttt{train\_unifolm\_wla.py:356}）——
\textbf{从训练迭代器里消费了一个 batch 并丢弃}。这不仅浪费数据，
还让 eval 的样本其实就是"下一步本该训练的那批"（轻微信息泄漏 + 训练步数据减少）。
调用频率是每 500 步一次，影响不大，但语义上是错的。应该用独立的 eval dataloader。

\subsection*{7.6 🟡 单进程运行会在训练结束时崩}

\texttt{main()} 末尾（\texttt{train\_unifolm\_wla.py:474-475}）无条件调用
\texttt{dist.barrier()} 和 \texttt{dist.destroy\_process\_group()}。
进程组未初始化时 \texttt{dist.barrier()} 会抛
\texttt{RuntimeError: Default process group has not been initialized}。
同理 \texttt{\_log\_metrics} 里的 \texttt{dist.get\_rank()}（\texttt{:281}）也是无条件调用。
\textbf{必须用 \texttt{torchrun} / \texttt{accelerate launch} 启动}，不能 \texttt{python train\_unifolm\_wla.py}。

\subsection*{7.7 🟡 \texttt{unitree.yaml} 的 \texttt{precollected\_stats\_path} 是相对 CWD 的裸相对路径}

\texttt{precollected\_stats\_path: "unifolm\_wla/dataloader/multi\_source\_dataset/stats"}（\texttt{unitree.yaml:23}）
经 \texttt{Path(...)} 解析时相对\textbf{当前工作目录}，而非 \texttt{data\_base} 或 YAML 所在目录。
\texttt{docs/train\_action\_expert.md:80} 要求在项目根目录执行，所以能工作——
但从别处启动就会 \texttt{FileNotFoundError}。建议改成相对 YAML 文件或加 \texttt{data\_base} 前缀。

\subsection*{7.8 🟡 numpy 2.x + \texttt{fastparquet} 等依赖的兼容性}

\texttt{pyproject.toml} 同时钉了 \texttt{numpy==2.2.6}、\texttt{pyarrow==21.0.0}、\texttt{fastparquet==2024.11.0}、
\texttt{torch==2.8.0}、\texttt{transformers==5.5.3}。注意 \texttt{transformers} 5.x 是大版本跳跃，
\texttt{QWen3.py:68} 用的是 \texttt{dtype=torch.bfloat16}（5.x 新 API，旧版是 \texttt{torch\_dtype}），
说明代码是配合 5.x 写的。\textbf{不要降级 transformers}。

\subsection*{7.9 🟡 collate → PIL → 再 tensor 的往返开销}

数据流是：\texttt{\_process\_images} 产出 tensor → \texttt{collate\_fn} stack 成 \texttt{(B,C,H,W)} →
\texttt{\_adapt\_multi\_source\_batch} 又 \texttt{TF.to\_pil\_image} 拆回 PIL → processor 再转回 tensor →
\texttt{\_build\_state\_tensor} 又把 state 从 list 重建 tensor。

对 \texttt{(B,3,336,448)} 的图像，这是一次无谓的 GPU→CPU(PIL)→CPU tensor 往返。
注释说明 role 顺序必须保持，但完全可以在 tensor 层面维护顺序、
只把 PIL 转换留给 processor。\textbf{这是可优化的性能热点}（每步都要做）。

\subsection*{7.10 🟡 \texttt{\_load\_stats} 要求必须提供 \texttt{precollected\_stats\_path}}

\texttt{single\_source\_dataset.py:268-273} 直接 \texttt{raise ValueError} 如果没设这个字段。
\texttt{peer\_data\_paths} 和 \texttt{merge\_left\_right\_ee\_stats} 两个字段在 \texttt{config.py} 里被解析、
在 \texttt{unitree.yaml} 里被设置，但\textbf{合并逻辑已被删除}（注释说"lived in the removed
multi\emph{source}dataset.py"）。所以这两个配置项现在是纯装饰，会误导使用者。
文档 \texttt{docs/robot\_action\_state\_processing.md} §12/§13 大篇幅讲了统计量合并，
但那是\textbf{离线统计脚本}的规范，不在本仓库内。

\subsection*{7.11 🟢 其他小问题}

\begin{itemize}
\item \texttt{baseframework.\_\_init\_\_(self, hf\_config=PretrainedConfig())} 用了可变默认参数
（\texttt{base\_framework.py:88}），所有实例共享同一个 config 对象
\item \texttt{share\_tools.py:99} \texttt{isinstance(x, OmegaConf.\_\_class\_\_)} 是个恒假的判断，注释也承认"typically not hit"
\item \texttt{trainer\_tools.py:137} \texttt{ReferenceError(...)} 只构造不抛异常，是个笔误（应为 \texttt{raise}）
\item \texttt{train\_unifolm\_wla.py:39} 的 \texttt{load\_fast\_tokenizer()} 定义了但从未调用（FAST tokenizer 的遗留）
\item \texttt{train\_unifolm\_wla.py:91} 的 rank 判断用了 \texttt{dist.is\_initialized() and dist.get\_rank()==0}，
而 \texttt{:281} 却直接 \texttt{dist.get\_rank()}，风格不一致
\end{itemize}

\hrule
\section*{8. 死代码 / 未接线清单}

这些先不用读，会浪费时间：

$\vert$ 文件 / 符号 $\vert$ 行数 $\vert$ 状态 $\vert$
$\vert$---$\vert$---$\vert$---$\vert$
$\vert$ \texttt{model/modules/action\_model/flow\_matching\_head/cross\_attention\_dit.py} $\vert$ 362 $\vert$ \textbf{整个文件零引用}，应是 QwenGR00T 时代 \texttt{LayerwiseFlowmatchingActionHead} 的残留 $\vert$
$\vert$ \texttt{dataloader/multi\_source\_dataset/fps\_utils.py} $\vert$ 83 $\vert$ \textbf{整个文件零引用}（见 §7.4） $\vert$
$\vert$ \texttt{model/tools.py: FrameworkTools} $\vert$ \textasciitilde{}90 $\vert$ 零引用（\texttt{base\_framework.py} 只 import 了名字不用）。其中 \texttt{unnormalize\_actions} 的 \texttt{gripper\_channel\_idx=6} 恰好对上新版 54 维布局的左夹爪槽位，纯属巧合 $\vert$
$\vert$ \texttt{model/tools.py: CrossAttention} $\vert$ \textasciitilde{}62 $\vert$ 零引用（早期 VGGT 融合的残留，forward 的 docstring 还写着 \texttt{vggt\_feature}） $\vert$
$\vert$ \texttt{model/tools.py: preprocess\_images} $\vert$ \textasciitilde{}90 $\vert$ 零引用（Prismatic/OpenVLA 的 518px 预处理遗留） $\vert$
$\vert$ \texttt{model/tools.py: auto\_get\_module\_keys} $\vert$ 22 $\vert$ 零引用 $\vert$
$\vert$ \texttt{model/tools.py: print\_freeze\_status} $\vert$ 33 $\vert$ 零引用 $\vert$
$\vert$ \texttt{share\_tools.py: read\_model\_config} $\vert$ 57 $\vert$ 零引用（被 \texttt{read\_mode\_config} 取代，后者读 YAML 而非 JSON） $\vert$
$\vert$ \texttt{share\_tools.py: populate\_layerwise\_dit\_cfg} $\vert$ 35 $\vert$ 零引用（为 \texttt{LayerwiseFlowmatchingActionHead} 服务） $\vert$
$\vert$ \texttt{share\_tools.py: state2str\_transform} / \texttt{add\_discretized\_state\_to\_instruction} $\vert$ \textasciitilde{}30 $\vert$ 零引用（π₀.5 风格的"状态离散化进 prompt"方案，未启用） $\vert$
$\vert$ \texttt{share\_tools.py: NamespaceWithGet} $\vert$ \textasciitilde{}40 $\vert$ 已定义但当前路径不再使用 $\vert$
$\vert$ \texttt{projector/robot\_state.py: add\_robot\_state\_projector} / \texttt{patch\_robot\_state\_forward} / \texttt{\_prepare\_robot\_state\_inputs} / \texttt{get\_llm\_model} / \texttt{get\_llm\_input\_embeddings} $\vert$ \textasciitilde{}120 $\vert$ 零引用。这些是 ms-swift 风格的"patch model.forward"方案；\texttt{QWen3.py} 走了更干净的 forward-hook 方案，只复用了 \texttt{RobotStateProjector} 和 \texttt{load\_robot\_state\_projector\_state\_dict} $\vert$
$\vert$ \texttt{trainer\_utils/trainer\_tools.py: setup\_optimizer\_and\_scheduler} $\vert$ 33 $\vert$ 零引用（\texttt{train\_unifolm\_wla.py} 里有一个同名但 \texttt{fused=True} 硬编码的版本） $\vert$
$\vert$ \texttt{trainer\_utils/trainer\_tools.py: eval\_qwenpi} / \texttt{compute\_grad\_angle\_with\_stats} / \texttt{pcgrad\_project} / \texttt{only\_main\_process} / \texttt{is\_main\_process} / \texttt{extract\_json\_from\_string} $\vert$ \textasciitilde{}120 $\vert$ 零引用（PCGrad 梯度投影、IoU 评测是早期研究实验残留） $\vert$
$\vert$ \texttt{base\_framework.py: compute\_loss} / \texttt{supports\_training\_tag} / \texttt{forward\_vlm} $\vert$ \textasciitilde{}66 $\vert$ 有一个"多 tag 统一 loss 接口"的设计（vla / vlm / world 三种数据），但 trainer 直接调 \texttt{model.forward()}，这套接口\textbf{从未被调用} $\vert$
$\vert$ \texttt{training/train\_unifolm\_wla.py: load\_fast\_tokenizer} $\vert$ 2 $\vert$ 零引用 $\vert$
$\vert$ \texttt{action\_mapping.py: map\_single\_left\_action} $\vert$ 17 $\vert$ 零引用 $\vert$
$\vert$ \texttt{config.py: ImageKeyConfig.resize} / \texttt{DatasetSourceConfig.image\_size} $\vert$ — $\vert$ 已标记 deprecated $\vert$
$\vert$ \texttt{DatasetSourceConfig.peer\_data\_paths} / \texttt{merge\_left\_right\_ee\_stats} $\vert$ — $\vert$ 被解析但合并逻辑已移除 $\vert$
$\vert$ collate 里的 \texttt{action\_norm\_offset} / \texttt{action\_norm\_scale} $\vert$ — $\vert$ 被计算、被 collate，\textbf{无消费者} $\vert$
$\vert$ \texttt{state\_unnorm} $\vert$ — $\vert$ 被产出但 \texttt{collate\_fn} 不收集，直接丢弃 $\vert$
$\vert$ \texttt{action\_model.action\_hidden\_dim} $\vert$ — $\vert$ 配置项被 \texttt{apply\_config\_compat} 补全，但模型代码不读它 $\vert$

\textbf{净结论}：真正参与一次训练前向/反向的代码，大约只有 \textbf{2500～3000 行}，
其余是移植残留、实验分支和预留接口。

\hrule
\section*{9. 上手路线图}

\subsection*{第一遍：抓主干（建议 2 小时）}

按这个顺序读，每一步都能跑起来或对上形状：

\begin{enumerate}
\item \texttt{README.md} + \texttt{docs/train\_action\_expert.md} —— 知道要干什么
\item \texttt{docs/robot\_action\_state\_processing.md} §1–§3 —— \textbf{动作/状态空间的定义源头}，先建立"54/60 维"的直觉
\item \texttt{action\_mapping.py} —— 把文档的表格变成代码里的 \texttt{slice}
\item \texttt{se3\_utils.py:162-180} —— 相对动作那 10 行
\item \texttt{single\_source\_dataset.py:451-530} —— \texttt{\_\_getitem\_\_} 的 8 个步骤
\item \texttt{dataloader.py:45-124} —— collate 输出的 batch 结构
\item \texttt{MMDiT\_ActionHeader.py:230-350} —— 训练 forward + 推理采样，\textbf{这是全仓库的心脏}
\item \texttt{QwenMMDiT.py:233-378} —— 组装与两条路径的差异
\item \texttt{train\_unifolm\_wla.py:308-416} —— 训练循环
\end{enumerate}

\subsection*{第二遍：抠细节（按兴趣选）}

\begin{itemize}
\item \textbf{想做数据侧}：\texttt{docs/robot\_action\_state\_processing.md} §4–§16 + \texttt{stats\_utils.py} + \texttt{lerobot\_wrapper.py}
\item \textbf{想做模型侧}：\texttt{mmdit.py} 的 \texttt{QwenDoubleStreamAttnProcessor2\_0} + \texttt{QwenImageTransformerBlock}，
对照 Qwen-Image 的原始实现看差异
\item \textbf{想做工程}：\texttt{overwatch.py} + \texttt{dataloader.py:\_CpuAcceleratorWorkerDataLoader} + \texttt{config\_tracker.py}
\item \textbf{想做多本体泛化}：\texttt{embodiment\_embedding} + FiLM + \texttt{DEFAULT\_EMBODIMENT\_TYPES}
\end{itemize}

\subsection*{第三遍：动手验证}

```python

\chapter*{1) 只测数据层，不需要模型权重}

from unifolm\emph{wla.dataloader.multi}source\emph{dataset.dataloader import create}training\emph{dataloader
dl, ds = create}training\emph{dataloader(".../unitree.yaml", batch}size=2, num\_workers=0)
b = next(iter(dl))
for k, v in b.items():
    print(k, v.shape if hasattr(v, "shape") else (type(v), len(v) if hasattr(v, "\textbf{len}") else v))

\chapter*{重点核对: action (2,30,54), action\_mask (2,54), state (2,60), images \{role:(2,3,336,448)\}}

\chapter*{2) 只测动作头，不需要 VLM}

import torch
from types import SimpleNamespace

\chapter*{造一个 fake full\_config，调用 MMDiTFlowmatchingActionHead.forward}

\chapter*{验证: vl\_embs (B,L,2048), actions (B,30,108), state (B,1,120) → loss 是标量}

\chapter*{3) 读 \textbf{getitem} 的中间量，核对 §5 的表格}

```

\hrule
\section*{10. 自测题}

能答出这些，说明你真懂了：

\begin{enumerate}
\item 为什么状态用 rotation-6D 而动作输出用 rotvec？两者变换函数分别在哪？
\item \texttt{[48:54]} 在动作向量和状态向量里分别是什么？为什么同一个数字区间含义不同？
\item \texttt{T\_rel = T\_curr⁻¹ · T\_future} 相比直接用绝对位姿，解决了什么跨本体问题？
\item 54 维掩码的作用是什么？它在 loss 里出现在哪一行？
\item \texttt{repeated\_diffusion\_steps=8} 省下了什么？代价是什么？
\item 训练时 \texttt{t} 的分布偏向哪里？为什么？\texttt{noise\_s} 起什么作用？
\item 推理时从 \texttt{t=0} 还是 \texttt{t=1} 出发？积分方向如何？
\item 为什么 DiT 联合注意力必须 \texttt{[action, text]} 而不能 \texttt{[text, action]}？给出两个理由。
\item RoPE 为什么把 action 锚在位置 0？
\item \texttt{embodiment\_embedding} 解决什么问题？为什么两个 FiLM 都 zero-init？
\item 为什么加载 Qwen3-VL 时必须保留 \texttt{input\_ids} 而不能只传 \texttt{inputs\_embeds}？
\item robot\_state 的两条注入通路分别是什么？为什么互斥？
\item \texttt{lamp}/\texttt{sync\_gradients} 守卫如果不加会怎样？
\item \texttt{AccessTrackedConfig} 为什么要 monkey patch OmegaConf 的 5 个函数？
\item DataLoader worker 里那 520 MiB 显存是怎么产生的？为什么不能用清空 \texttt{CUDA\_VISIBLE\_DEVICES} 来修？
\item \texttt{read\_mode\_config} 期望 checkpoint 目录里有哪些文件？
\item 仓库自带训练 YAML 为什么开箱跑不通？
\item \texttt{pred[:, -H:]} 里的负号为什么是必要的？
\end{enumerate}

\hrule
\section*{附：一页速查}

\texttt{
统一动作      54 维  (动作输出用 xyz+rotvec，EE/底盘位姿是相对当前状态的 SE(3))
统一状态      60 维  (EE 用 xyz+rot6d；[41:47] 是重力向量+角速度，不是位姿)
chunk 长度    = source\_fps（当前 30），chunk\_size 配置项未被使用
归一化        minmax\_q(动作/状态/夹爪) + zscore(相对位姿)，global\_* 优先
VLM           Qwen3-VL-4B，hidden=2048，只取 last\_hidden\_states 当编码器
条件注入      cross-attention（VLM hidden → RMSNorm → Linear(2048→768)）+ 可选 state token
DiT           DiT-B: 12 头 × 64 = inner 768；16 层双流联合注意力；output 1024
动作 token    ActionEncoder(108→768) + pos\_embed(1024,768) + FiLM\_pre
损失          掩码加权 MSE(pred\_velocity, action - noise)，t \textasciitilde{} 0.999*(1-Beta(1.5,1))
推理          4 步 Euler，dt=0.25，每步乘掩码
训练          Accelerate + DeepSpeed ZeRO-2 + bf16，action\_horizon=30，repeat=8
分层 LR       base 1e-4 / VLM 1e-5 / action\_model 1e-4，cosine\_with\_min\_lr
}

\end{document}
